\section{Splitting a secret}
\label{sec:splitting-secrets}

Splitting is how a standalone secret is escrowed, and how a file's data key
is protected --- both use the same tree machinery.

\begin{noteBox}
The live SQLite tree \emph{is} the spec for whoever holds that store. An
operator or publisher may keep the full org tree; a personal store should
hold only the nodes that person needs for signing and access control: their
lineage, their descendants, siblings of their node, and the fixpoint of
\emph{established} bridge peers (the peer and the peer's ancestors, not the
peer's unrelated siblings). Pushing from a personal store merges that
subgraph into the relay; it does not replace the canonical document. There
is no JSON file to author by hand --- \code{split}, \code{bind}, \code{add},
\code{revoke}, \code{bridge}, and \code{access quorum --state 0 --leaf}
write the tree in place. \code{tree --output} writes a snapshot of whatever
is stored now; \flag{--tree-spec FILE} remains only for a nested one-shot
tree.
\end{noteBox}

Labels must be unique within a tree. \code{tree <id> --node A B} prints
that set's lowest common ancestor; \code{reconstruct --node A B} starts
recovery at that ancestor. \code{tree} with no id lists every stored spec.

\subsection{A two-department example}

\code{split --source master.pub} escrows that key file (hex, PEM, or
OpenSSH). \flag{--leaf} builds a one-level tree and binds every sibling
pair, so \code{M.S <-> M.A} is recorded even after a later refresh.
Reconstruct with the holders' key files and \flag{--output}. A 2-of-2
department tree (\code{M} = \file{master.pub}, \code{M.S} =
\file{SoftwareDepartment.pub}, \code{M.A} =
\file{AccountingDepartment.pub}) reassembles the master from the two
department keys:

\begin{lstlisting}
keyquorum split --label master --threshold 2 \
  --leaf M.S=SoftwareDepartment.pub --leaf M.A=AccountingDepartment.pub \
  --source master.pub --generate-keys --register
keyquorum reconstruct <key-id> \
  --share-file SoftwareDepartment.pub --share-file AccountingDepartment.pub \
  --output master.pub
keyquorum tree
keyquorum tree <key-id> --output org.json
\end{lstlisting}

\flag{--generate-keys} creates each leaf's \file{.pub} / \file{.key} pair;
\flag{--register} records those public keys. Dotted leaf labels
(\code{M.S}, \code{M.A}) infer the root node \code{M}. \flag{--bind
M.S=M.A} adds an extra pairing; \flag{--leaf} already binds sibling leaves.
A \file{.pub} share file uses the sibling \file{.key} (or a private-key
prompt) to unwrap that department's sealed share --- both departments are
required because \code{M}'s threshold is 2.

\subsection{Growing and rebinding a tree}

Later commands keep changing the same tree. Pairings are stored by node id,
so they survive a secret refresh, a new leaf, or a leaf moving to a new
public key:

\begin{lstlisting}
# Pair two nodes (whitelist both ways + establish the link)
keyquorum bind <key-id> --node M.S --peer M.A

# Move M.S onto a new token; node id -- and the bind -- stay put
keyquorum bind <key-id> --node M.S \
  --public-key-file NewSoftware.pub --share-file SoftwareDepartment.key \
  --register

# Grow M to 2-of-3; survivor node ids (and M.S <-> M.A) stay put
keyquorum add <key-id> --parent M --node M.F \
  --public-key-file FinanceDepartment.pub \
  --share-file SoftwareDepartment.pub --share-file AccountingDepartment.pub \
  --generate-keys --register
\end{lstlisting}

\subsection{Bridges (whitelist and pairing)}

\begin{lstlisting}
keyquorum tree 1 --node alice bob
keyquorum reconstruct 1 --node alice bob --share-file alice.pub \
  --share-file bob.key
keyquorum bridge allow 1 --node alice --peer it
keyquorum bridge add 1 --from alice --to it
keyquorum bridge list 1
keyquorum bridge remove 1 --from alice --to it
keyquorum bridge deny 1 --node alice --peer it
\end{lstlisting}

\code{bind --peer} is the usual way to stand up \code{A <-> B}.
\code{bridge allow} / \code{deny} change the whitelist only. \code{bridge
add} / \code{remove} stand up or tear down an established pairing --- add
requires a whitelist hit on either side; deny also drops any pairing.

\begin{noteBox}
This bridge/whitelist pairing is unrelated to a \emph{private sign bridge}
(Section~\ref{sec:private-bridges}), which is a separate, N-person
co-signing group.
\end{noteBox}

\subsection{Publishing and fetching topology}

The relay stores the \textbf{full} public tree as a JSON document. Sending
data (\code{relay push}, including private-bridge envelopes) merges the
sender's public topology into that document; nodes the sender does not hold
stay in place, so a personal subgraph cannot replace canonical context.
Fetching data (\code{relay pull}) returns the slice this device's
encryption fingerprint is allowed to see, and the CLI merges it into local
SQLite before importing envelopes. A later push that adds a bridge such as
\code{M.S.2 <-> M.A.1} expands the next pull for \code{M.S.2} to include
\code{M.A.1} as topology-only (no sealed share).

\code{tree fetch} refreshes topology without downloading envelopes.
\code{tree publish} remains if you need to replace the document without an
envelope.

\begin{lstlisting}
# Person: envelopes plus the slice this pull key is allowed to see
keyquorum --db alice.sqlite relay pull --import --share-file alice.key

# Topology only (no envelopes), including first fetch onto an empty store:
keyquorum tree fetch <key-id>
keyquorum tree fetch --label master
\end{lstlisting}
